\chapter{Bloch-Periodic Finite Differences on Nonorthogonal Cells}
\index{finite difference!Bloch-periodic|textbf}
\index{Laplacian!nonorthogonal periodic cell|textbf}
\index{Bloch boundary condition|textbf}

\label{ch:training-bloch-finite-differences}

\chapterstatus{pedagogical derivation and verified software-capability
description. The maintained operator is a centered second-order geometric scalar
Laplacian. This chapter reports no new M4 calculation, material Hamiltonian,
scientific validation, or uncertainty quantification.}

\section{Learning goals}

After working through this chapter, a reader should be able to

\begin{enumerate}
  \item derive the Cartesian scalar Laplacian in fractional coordinates;
  \item distinguish metric-induced mixed derivatives from physical
        effective-mass anisotropy;
  \item construct pure and mixed centered second-order stencils;
  \item apply Bloch phases when a stencil crosses one or several periodic
        seams;
  \item construct plane-wave and finite-difference representations of a
        one-dimensional periodic cosine fiber;
  \item distinguish the geometric Laplacian from kinetic-energy and potential
        operators; and
  \item design software and numerical checks without turning an expected trend
        into a scientific conclusion.
\end{enumerate}

Chapter~\ref{ch:training-periodic-geometry} supplies the direct basis $A$, Gram
metric $G=A^{\mathsf T}A$, inverse metric $G^{-1}$, and explicit route from an
ambient $3\times2$ basis to an intrinsic $2\times2$ basis. Standard
finite-difference background may be found in Strikwerda~\cite{strikwerda2004};
the evidence classifications used here follow
Chapter~\ref{ch:verification-and-validation}.

\section{Why a familiar stencil needs a full lecture}

Most readers have seen the three-point approximation to a second derivative.
It is therefore reasonable to ask why we should spend an entire lecture on a
Laplacian. The answer is that the one-dimensional stencil is only the local
algebraic ingredient. A periodic multidimensional operator also commits us to a
coordinate system, a metric, a tensor ordering, a boundary phase convention,
a matrix ordering, a physical unit, and a sign convention. Each choice is
simple when isolated. Their composition is where silent errors appear.

On a square grid with zero twist, several mistakes can hide. If one treats
fractional coordinates as Cartesian coordinates, the square unit cell may make
the error invisible. If one omits mixed derivatives, an orthogonal metric makes
the omitted coefficients zero anyway. If one mishandles simultaneous seam
crossings, an axial stencil never tests the corner case. A nonorthogonal cell
with nonzero twist is therefore not needless generality: it exposes assumptions
that an easy special case allows us to forget.

The lecture will proceed in the same order one would use at a blackboard. We
first derive the continuous operator in reduced coordinates. We then replace
each derivative by a centered stencil. Only after the interior stencil is clear
do we discuss periodic seams. Finally, we attach units and explain how another
model layer turns a geometric Laplacian into a Hamiltonian. This order keeps the
source of each coefficient visible.

A useful question to ask at every stage is: ``What has been added since the
previous line?'' Moving from $G^{-1}$ to $\Delta$ adds differential structure.
Moving from $\Delta$ to $\Delta_h$ adds a discretization. Moving from an
interior stencil to a Bloch matrix adds a boundary representation. Moving from
$\Delta_h$ to $H_h$ adds physical scaling and a potential. None of these steps
is merely a change of notation.

\section{From Cartesian derivatives to fractional derivatives}
\index{coordinate transformation!Laplacian}

Let
\begin{equation}
  \mathbf r=A\mathbf s,
  \qquad
  G=A^{\mathsf T}A,
  \label{eq:training-fd-coordinate-map}
\end{equation}
where the primitive basis is constant over the cell. For a scalar function
$f(\mathbf r)=\widetilde f(\mathbf s)$, the gradient transformation gives a
constant-coefficient Laplacian
\begin{equation}
  \nabla_{\mathbf r}^2 f
  =
  \sum_{i,j=1}^{d}
  g^{ij}
  \frac{\partial^2\widetilde f}
       {\partial s_i\partial s_j},
  \qquad
  [g^{ij}]=G^{-1}.
  \label{eq:training-metric-laplacian}
\end{equation}
Since $G^{-1}$ is symmetric, the same expression can be written
\begin{equation}
  \nabla_{\mathbf r}^2
  =
  \sum_{i=1}^{d}g^{ii}\partial_{s_i}^2
  +2\sum_{i<j}g^{ij}\partial_{s_i}\partial_{s_j}.
  \label{eq:training-metric-laplacian-symmetric}
\end{equation}

For an orthogonal cell, $g^{ij}=0$ when $i\neq j$, and the mixed derivatives
vanish. For a skew cell, the mixed terms are required to represent the ordinary
isotropic Cartesian Laplacian. Omitting them changes the mathematical
operator.

\subsection{The chain rule one index at a time}

The compact expression in Equation~\eqref{eq:training-metric-laplacian} is easy
to quote and just as easy to misremember. We therefore derive it. Let
$\nabla_{\mathbf s}\widetilde f$ be the column of fractional derivatives. A
small fractional displacement changes the function by
\begin{equation}
  d\widetilde f
  =\left(\nabla_{\mathbf s}\widetilde f\right)^{\mathsf T}
   \mathrm{d}\mathbf s.
\end{equation}
The same change, described in the tangent Cartesian plane, is
\begin{equation}
  df
  =\left(\nabla_{\parallel}f\right)^{\mathsf T}\mathrm{d}\mathbf r
  =\left(\nabla_{\parallel}f\right)^{\mathsf T}A\,\mathrm{d}\mathbf s.
\end{equation}
Equality for every $\mathrm{d}\mathbf s$ requires
\begin{equation}
  A^{\mathsf T}\nabla_{\parallel}f
  =\nabla_{\mathbf s}\widetilde f.
  \label{eq:training-gradient-pullback}
\end{equation}
The in-plane solution is
\begin{equation}
  \nabla_{\parallel}f
  =AG^{-1}\nabla_{\mathbf s}\widetilde f.
  \label{eq:training-gradient-pushforward}
\end{equation}
Substitution confirms the result because
$A^{\mathsf T}AG^{-1}=I_d$.

The basis is constant, so differentiating once more does not produce derivatives
of $A$ or $G$. Taking the tangent divergence leaves one inverse-metric
coefficient for each pair of fractional derivatives, giving
Equation~\eqref{eq:training-metric-laplacian}. For a curvilinear or spatially
varying basis, this shortcut would be invalid; metric determinants and
coefficient derivatives would enter. The present implementation is explicitly
a constant primitive-cell construction.

This derivation is also valid for an embedded plane because
$AG^{-1}\nabla_{\mathbf s}\widetilde f$ lies in the column space of $A$. We are
not taking a three-dimensional Laplacian of a function whose normal dependence
has been left unspecified. We are taking the intrinsic scalar Laplacian of a
flat plane represented in ambient coordinates.

\begin{remark}[Geometry is not constitutive anisotropy]
Off-diagonal entries of $G^{-1}$ arise because the selected coordinate axes are
skew. An anisotropic effective-mass tensor describes a physical constitutive
model. Although both can produce cross derivatives after a coordinate choice,
they belong to different layers and must not be identified merely because
their matrices have similar shapes.
\end{remark}

\section{Endpoint-excluded periodic grids}
\index{grid!endpoint-excluded periodic}

For axis $i$, select $N_i$ points
\begin{equation}
  s_{i,n_i}=\frac{n_i}{N_i},
  \qquad
  n_i=0,1,\ldots,N_i-1,
  \qquad
  h_i=\frac{1}{N_i}.
  \label{eq:training-periodic-grid}
\end{equation}
The endpoint $s_i=1$ is excluded because it is equivalent to $s_i=0$ up to the
boundary condition. Including both endpoints would duplicate one physical
location and complicate the seam rule.

A tensor-product grid also needs a declared linear ordering. The maintained
implementation uses C order: the last fractional axis varies fastest. Ordering
is not a cosmetic implementation detail because it determines which matrix row
and column correspond to each grid point.

\section{Centered pure derivatives}

The centered second derivative along fractional axis $i$ is
\begin{equation}
  (D_{ii}f)_{\mathbf n}
  =
  \frac{
    f_{\mathbf n+\mathbf e_i}
    -2f_{\mathbf n}
    +f_{\mathbf n-\mathbf e_i}}
  {h_i^2}.
  \label{eq:training-pure-second-difference}
\end{equation}
For a smooth function away from representation boundaries, this is a
second-order approximation to $\partial_{s_i}^2f$. The corresponding axial
matrix entries are real before seam phases are applied.

The pure-derivative contribution to the discrete metric Laplacian is
\begin{equation}
  \sum_i g^{ii}D_{ii}.
  \label{eq:training-pure-metric-stencil}
\end{equation}
Because $G^{-1}$ is positive definite, each $g^{ii}$ is positive. With the
sign convention in Equation~\eqref{eq:training-pure-second-difference}, the
periodic discrete Laplacian has nonpositive eigenvalues.

\section{Centered mixed derivatives}
\index{finite difference!mixed derivative}

For $i\neq j$, the centered mixed derivative uses four corner points:
\begin{equation}
  (D_{ij}f)_{\mathbf n}
  =
  \frac{
    f_{\mathbf n+\mathbf e_i+\mathbf e_j}
    -f_{\mathbf n+\mathbf e_i-\mathbf e_j}
    -f_{\mathbf n-\mathbf e_i+\mathbf e_j}
    +f_{\mathbf n-\mathbf e_i-\mathbf e_j}}
  {4h_i h_j}.
  \label{eq:training-mixed-second-difference}
\end{equation}
The complete centered operator is
\begin{equation}
  \Delta_h
  =
  \sum_i g^{ii}D_{ii}
  +2\sum_{i<j}g^{ij}D_{ij}.
  \label{eq:training-discrete-metric-laplacian}
\end{equation}
For each pair $i<j$, the factor of two comes from combining the $(i,j)$ and
$(j,i)$ terms in Equation~\eqref{eq:training-metric-laplacian}.

In two dimensions, the pure stencil touches the center and four axial
neighbors; a nonzero mixed coefficient adds four diagonal neighbors. In three
dimensions, three axis pairs contribute corresponding corner displacements.
The represented operator remains sparse.

\subsection{Reading the stencil for a sixty-degree cell}

Return to the sixty-degree cell from
Equation~\eqref{eq:training-sixty-degree-basis}. Its inverse metric is
\begin{equation}
  G^{-1}=\frac{1}{a^2}
  \begin{bmatrix}
    4/3 & -2/3\\
    -2/3 & 4/3
  \end{bmatrix}.
\end{equation}
The continuous Laplacian is therefore
\begin{equation}
  \nabla^2
  =\frac{4}{3a^2}
  \left(
    \partial_{s_1}^2
    +\partial_{s_2}^2
    -\partial_{s_1}\partial_{s_2}
  \right).
  \label{eq:training-sixty-degree-laplacian}
\end{equation}
The minus sign on the mixed derivative is not a negative kinetic energy and not
an instability. It is the coordinate correction required by the oblique axes.
The full quadratic symbol remains nonnegative for $-\nabla^2$ because the
inverse metric is positive definite.

Suppose, for simplicity, that $h_1=h_2=h$. The center receives
$-16/(3a^2h^2)$. Each of the four axial neighbors receives
$4/(3a^2h^2)$. The four corner coefficients alternate in sign: same-sign
corner displacements receive $-1/(3a^2h^2)$, while opposite-sign corner
displacements receive $+1/(3a^2h^2)$. Adding every coefficient in a row gives
zero before a nonzero boundary twist is applied. That row-sum check is the
discrete statement that a constant periodic function has zero Laplacian.

It is worth tracing the signs directly from
Equation~\eqref{eq:training-mixed-second-difference}. The mixed derivative has
positive same-sign corners and negative opposite-sign corners. Multiplication
by the negative $g^{12}$ reverses those signs. This two-step explanation is
safer than memorizing a nine-point stencil whose signs depend on the selected
cell.

\section{Bloch-periodic seams}
\index{Bloch phase!quotient seam}
\index{boundary twist}

Let the reduced twist $\boldsymbol\kappa$ be measured in turns along the direct
primitive axes. The maintained quotient-seam convention is
\begin{equation}
  \psi(\mathbf s+\mathbf q)
  =
  \exp\!\left(2\pi i\,\boldsymbol\kappa\cdot\mathbf q\right)
  \psi(\mathbf s),
  \qquad
  \mathbf q\in\mathbb Z^d.
  \label{eq:training-bloch-seam}
\end{equation}
Suppose a stencil displacement carries a raw target outside the stored grid.
Reduce it as
\begin{equation}
  \mathbf n+\boldsymbol\delta
  =\overline{\mathbf n}+N\mathbf q,
  \label{eq:training-grid-quotient}
\end{equation}
where $\overline{\mathbf n}$ lies in the stored index range and $\mathbf q$
counts signed primitive-cell crossings. The matrix contribution receives the
phase
\begin{equation}
  \exp\!\left(2\pi i\,\boldsymbol\kappa\cdot\mathbf q\right)
  =
  \prod_a\exp\!\left(2\pi i\kappa_a q_a\right).
  \label{eq:training-seam-product}
\end{equation}

The product form matters for a mixed-derivative corner that crosses two seams
at once. Applying only one axis phase produces the wrong operator. Integer
shifts of the twist define the same quotient-seam phases, but the complete
unreduced lift and its quotient reduction can still be retained as provenance.

\subsection{Walking one corner across the seam}

Take a $5\times4$ grid. The stored indices are
$n_1\in\{0,1,2,3,4\}$ and $n_2\in\{0,1,2,3\}$. Begin at the upper-right stored
point $(4,3)$ and apply the corner displacement $(+1,+1)$. The raw target is
$(5,4)$. Componentwise Euclidean division gives the wrapped point $(0,0)$ and
the quotient
\begin{equation}
  \mathbf q=(1,1).
\end{equation}
The associated matrix entry receives
\begin{equation}
  \exp\!\left[2\pi i(\kappa_1+\kappa_2)\right].
  \label{eq:training-double-wrap-phase}
\end{equation}
Both seam crossings matter.

Now begin at $(4,0)$ and apply $(+1,-1)$. The wrapped point is $(0,3)$, but the
quotient is $(1,-1)$. Its phase is
\begin{equation}
  \exp\!\left[2\pi i(\kappa_1-\kappa_2)\right].
\end{equation}
The negative quotient records that the second coordinate crossed the lower
seam. Using an unsigned ``did wrap'' flag would lose this distinction.

Finally consider the reverse matrix connection. It carries the opposite
quotient and therefore the complex-conjugate phase. Paired with the real
symmetric interior coefficient, this is the local reason the assembled matrix
can be Hermitian. The global Hermiticity test should still be performed; the
hand argument tells us what the test is checking.

\subsection{Hermiticity check}

For a real symmetric metric and a consistent seam convention, the represented
matrix should satisfy
\begin{equation}
  L_h^{\dagger}=L_h.
  \label{eq:training-fd-hermiticity}
\end{equation}
This is an important software invariant. It is not, by itself, a convergence
or scientific-validation result. Many incorrect Hermitian matrices can still
be assembled.

\section{From the Laplacian to a Hamiltonian}
\index{kinetic energy!scalar Laplacian distinction}

The geometric result constructed above represents $\nabla^2$. For an isotropic
scalar mass $m$, a familiar kinetic operator is
\begin{equation}
  \hat T
  =-\frac{\hbar^2}{2m}\nabla^2.
  \label{eq:training-isotropic-kinetic}
\end{equation}
A finite Hamiltonian may then have the form
\begin{equation}
  H_h=-\alpha L_h+V_h,
  \qquad
  \alpha=\frac{\hbar^2}{2m},
  \label{eq:training-fd-hamiltonian-composition}
\end{equation}
where $V_h$ is a separately defined potential representation. The basis,
units, sampling convention, and energy reference of both terms must agree
before addition.

An anisotropic effective-mass model instead involves a tensor such as
\begin{equation}
  \hat T_M
  =-\frac{\hbar^2}{2}
  \nabla\!\cdot M^{-1}\nabla,
  \label{eq:training-anisotropic-kinetic}
\end{equation}
subject to the model's ordering and regularity assumptions. That tensor is not
supplied by the scalar geometric Laplacian. It requires a separate physical
contract and, for spatially varying coefficients, a separate discretization
analysis.

\section{A one-dimensional cosine-fiber bridge}
\index{cosine potential!plane-wave and finite-difference representations}

Before returning to the maintained multidimensional constructor, it is useful
to specialize the preceding rules to the one-dimensional parent used later in
the isolated-band reduction. A Bloch fiber was defined in
Equation~\eqref{eq:training-geometry-fiber-definition} as the restriction of the
periodic parent to one fixed quasimomentum sector, represented here on
cell-periodic functions. Let
\begin{equation}
  \widehat H
  =-\frac{\hbar^2}{2m}\frac{\mathrm d^2}{\mathrm d x^2}
   +V_0\cos(Gx),
  \qquad
  G=\frac{2\pi}{a}.
  \label{eq:training-fd-one-dimensional-parent}
\end{equation}
There are two common one-cell representations.

\subsection{Cell-periodic plane waves}

Write $\psi_k(x)=e^{ikx}u_k(x)$ with $u_k(x+a)=u_k(x)$. In the periodic basis
\begin{equation}
  \phi_n(x)=a^{-1/2}e^{inGx},
  \qquad n\in\mathbb Z,
\end{equation}
the cell-periodic fiber has entries
\begin{equation}
  H^{\mathrm{PW}}_{nm}(k)
  =\frac{\hbar^2(k+nG)^2}{2m}\delta_{nm}
   +\frac{V_0}{2}
    \left(\delta_{n,m+1}+\delta_{n,m-1}\right).
  \label{eq:training-fd-one-dimensional-plane-wave}
\end{equation}
The kinetic term is diagonal because each basis function is a derivative
eigenfunction. The cosine contains only the harmonics $+G$ and $-G$, so it
connects reciprocal indices separated by one. This reciprocal-index sparsity
is a statement about the potential's Fourier content, not yet a short-range
real-space hopping model.

A cutoff $P$ retains $n=-P,\ldots,P$. The resulting matrix dimension is
$2P+1$, and changing $P$ changes the represented space. Low eigenvalues may
stabilize rapidly while states near the reciprocal cutoff remain poorly
represented.

\subsection{A grid with a Bloch seam}

Alternatively, store the original wavefunction on
\begin{equation}
  x_j=\frac{ja}{N_x},
  \qquad j=0,\ldots,N_x-1,
\end{equation}
and impose $\psi(x+a)=e^{ika}\psi(x)$. If
\begin{equation}
  c_h=\frac{\hbar^2}{2m(a/N_x)^2},
\end{equation}
the centered kinetic matrix has diagonal entries $2c_h$, adjacent entries
$-c_h$, and seam entries
\begin{equation}
  T_{0,N_x-1}(k)=-c_h e^{-ika},
  \qquad
  T_{N_x-1,0}(k)=-c_h e^{ika}.
  \label{eq:training-fd-one-dimensional-seam}
\end{equation}
The potential is diagonal with values $V_0\cos(Gx_j)$. The two seam entries are
complex conjugates, which is the one-dimensional instance of the quotient rule
in Equation~\eqref{eq:training-seam-product}.

\subsection{What can be compared directly}

The two finite matrices approximate the same declared continuum fibers but use
different bases and generally different dimensions. Their ordered low
eigenvalues can be compared at common $k$ values. The matrices themselves
cannot be subtracted merely because both are called Hamiltonians. An operator
residual requires an explicit map into common coordinates.

A high plane-wave cutoff can serve as a declared finite represented reference
for a bounded refinement study. Such a reference is not the exact continuum
spectrum. Plane-wave cutoff error and finite-difference grid error also remain
separate channels even when both are measured against that reference.
Chapter~\ref{ch:training-isolated-band-fourier-reduction} continues from these
fiber eigensystems to an isolated scalar band and its finite hopping
representation.

\section{What the maintained implementation owns}
\index{PhysKit!Bloch-periodic Laplacian}

The pinned PhysKit implementation owns

\begin{itemize}
  \item endpoint-excluded tensor grids in two or three dimensions;
  \item intrinsic nonorthogonal direct metrics;
  \item centered second-order pure and mixed derivatives;
  \item quotient-seam Bloch phases, including simultaneous multi-axis wraps;
  \item C-order sparse assembly;
  \item explicit physical or unitless basis units; and
  \item an immutable canonical complex CSR representation correlated with its
        request and twist context.
\end{itemize}

For an ambient $3\times2$ plane, the explicit adapter from
Chapter~\ref{ch:training-periodic-geometry} supplies the intrinsic metric while
retaining the ambient basis and maps. The Laplacian constructor does not need
to pretend that the rectangular basis is square.

The implementation intentionally does not own kinetic-energy scaling,
effective-mass tensors, potential sampling, material parameters, model
selection, campaign thresholds, or scientific acceptance. Those belong to the
consumer that constructs a Hamiltonian and states the question being tested.

The next program deliberately uses the maintained public objects rather than
reproducing their implementation. The example remains small enough that every
construction step is visible.

\begin{pythonexample}{constructing a twisted scalar Laplacian}
import numpy as np
from projectkoios.physkit.periodic.finite_difference.grid import (
    PeriodicFiniteDifferenceGrid,
)
from projectkoios.physkit.periodic.finite_difference.laplacian import (
    BlochPeriodicLaplacianConstructionRequest,
    BlochPeriodicLaplacianConstructor,
)
from projectkoios.physkit.periodic.lattice.boundary_phases import (
    BoundaryTwistLift,
)
from projectkoios.physkit.periodic.lattice.embedding import (
    EmbeddedDirectLattice2D,
    EmbeddedDirectLattice2DAdaptationRequest,
    EmbeddedDirectLattice2DAdapter,
)
from projectkoios.physkit.periodic.lattice.finite_domain import (
    FinitePeriodicDomain,
    LatticeDimension,
)
from projectkoios.physkit.units import MatrixQuantity, PhysicalUnit

# Define a flat two-dimensional lattice in ambient three-space.
unit = PhysicalUnit("angstrom")
ambient_basis = np.array(
    [[1.0, 0.25], [0.0, 1.2], [0.5, -0.1]],
    dtype=np.float64,
)
embedded = EmbeddedDirectLattice2D(
    ambient_basis=MatrixQuantity(ambient_basis, unit)
)

# Adapt ambient geometry to the intrinsic metric used by the stencil.
adaptation = EmbeddedDirectLattice2DAdapter().action(
    request=EmbeddedDirectLattice2DAdaptationRequest(
        embedded_lattice=embedded
    )
)
grid = PeriodicFiniteDifferenceGrid(
    metric=adaptation.intrinsic_metric,
    domain=FinitePeriodicDomain(LatticeDimension.TWO, (4, 5)),
    direct_basis_unit=unit,
)

# The twist is measured in turns along the two primitive axes.
request = BlochPeriodicLaplacianConstructionRequest(
    grid=grid,
    twist=BoundaryTwistLift(
        LatticeDimension.TWO,
        (0.125, -0.25),
    ),
)
result = BlochPeriodicLaplacianConstructor().action(request=request)
L = result.represented_laplacian.to_csr()

# Preserve sparsity while checking the exact construction invariant.
hermitian_residual = L - L.conjugate().transpose()
assert hermitian_residual.nnz == 0
print("shape:", L.shape)
print("unit:", result.represented_laplacian.unit)
\end{pythonexample}

\documentcallout{Notebook to do}{Create
\path{examples/tutorials/research-monograph/foundations/bloch_periodic_finite_differences.ipynb}
with explicit imports, explanatory inline comments, a zero-twist null-mode
check, and a sparse-preserving nonzero-twist Hermiticity check. The notebook
must not densify the represented Laplacian and must not imply kinetic scaling
or material physics.}

\section{A Fourier-mode diagnostic}
\index{Fourier mode!finite-difference symbol}

A periodic constant-coefficient stencil should be understood in Fourier space
as well as in matrix form. Consider a reduced-coordinate mode
\begin{equation}
  f_{\boldsymbol\nu}(\mathbf s)
  =\exp\!\left(2\pi i\,\boldsymbol\nu\cdot\mathbf s\right),
  \qquad
  \boldsymbol\nu=\mathbf m+\boldsymbol\kappa,
  \label{eq:training-discrete-bloch-mode}
\end{equation}
where $\mathbf m$ is an integer mode index and $\boldsymbol\kappa$ supplies the
Bloch twist. Translating by one primitive cell multiplies this mode by the
phase in Equation~\eqref{eq:training-bloch-seam}.

Applying the pure stencil gives the eigenvalue
\begin{equation}
  \widehat D_{ii}(\boldsymbol\nu)
  =-\frac{4}{h_i^2}
  \sin^2\!\left(\pi\nu_i h_i\right).
  \label{eq:training-pure-difference-symbol}
\end{equation}
Applying the mixed stencil gives
\begin{equation}
  \widehat D_{ij}(\boldsymbol\nu)
  =-\frac{
    \sin\!\left(2\pi\nu_i h_i\right)
    \sin\!\left(2\pi\nu_j h_j\right)}
  {h_i h_j}.
  \label{eq:training-mixed-difference-symbol}
\end{equation}
The full discrete symbol follows by inserting these expressions into
Equation~\eqref{eq:training-discrete-metric-laplacian}.

This is an especially useful numerical oracle because it is derived without
assembling the sparse matrix row by row. On a small grid, one can construct the
mode vector, apply the matrix, and compare the result with the analytical
stencil symbol. A discrepancy can reveal a wrong flattening order, a missing
mixed term, or an incorrect seam phase. Agreement remains bounded numerical
verification of this operator; it does not validate any material model.

In the small-spacing limit,
$\sin(\pi\nu_i h_i)\sim\pi\nu_i h_i$, so the pure symbol approaches
$-(2\pi\nu_i)^2$. The mixed symbol similarly approaches
$-(2\pi\nu_i)(2\pi\nu_j)$. This recovers the continuous metric symbol and makes
the expected second-order truncation behavior visible.

\section{A staged computational laboratory}
\index{laboratory!nonorthogonal finite difference}

A useful student exercise should separate contracts rather than beginning with
a large material calculation. The point of the laboratory is not to obtain one
impressive matrix. It is to make each layer fail in an interpretable way before
the layers are composed.

For every stage, record the request, the expected invariant, the observed
quantity, the acceptance rule, and the interpretation. If a stage fails, stop
and diagnose it rather than carrying a suspect operator into the next stage.
This is slower than writing the final Hamiltonian immediately and much faster
than debugging an unexplained band-structure discrepancy at the end.

\subsection{Stage 1: orthogonal null mode}

Use zero twist and an orthogonal unit cell. Construct the discrete scalar
Laplacian and verify that the constant vector is a null vector up to the
declared numerical tolerance. Check matrix shape, ordering, units, and
Hermiticity separately.

\subsection{Stage 2: one-dimensional symbol inside a tensor grid}

Choose a function that varies along only one fractional axis. Compare the
observed eigenvalue with the centered-difference symbol
\begin{equation}
  \lambda_h(\xi)
  =-\frac{4}{h^2}\sin^2\!\left(\frac{\xi h}{2}\right).
  \label{eq:training-centered-symbol}
\end{equation}
This isolates the pure derivative before introducing mixed terms.

\subsection{Stage 3: skew intrinsic cell}

Choose a full-rank skew $2\times2$ basis. Compare the represented action on a
fractional Fourier mode with the metric-weighted discrete symbol. Repeat with
the off-diagonal metric entry artificially omitted and explain why the latter
operator is not the same Cartesian Laplacian.

\subsection{Stage 4: ambient rotation}

Embed the same plane with a $3\times2$ basis, rotate it rigidly in ambient
three-space, and adapt both bases to the intrinsic route. The geometric scalar
Laplacians should agree because the induced metrics agree. The stored ambient
frames and projectors should rotate rather than disappear.

\subsection{Stage 5: seam and corner phases}

Use a nonzero two-component twist and inspect a mixed-stencil point that wraps
both axes. Verify that its phase equals the product in
Equation~\eqref{eq:training-seam-product}. Then check Hermiticity of the full
matrix.

\subsection{Stage 6: refinement without overclaiming}

Apply the operator to a smooth manufactured Bloch mode over a sequence of
grids. Estimate the observed order over the tested sequence. Report this as
numerical verification of the discretization for the manufactured case, not as
material validation or uncertainty quantification.

\subsection{What a laboratory report should say}

A strong report does not merely state that the code ``worked.'' It says, for
example, that the matrix shape and unit matched the declared grid contract;
that the maximum absolute Hermiticity residual was evaluated; that an
independently derived Fourier symbol agreed over the tested modes; and that a
manufactured smooth mode displayed a stated observed refinement rate over a
listed grid sequence. It also says what those observations do not establish.

If physical units are used, the report should show the dimensional path. The
basis has length, the inverse metric has inverse length squared, fractional
spacings are dimensionless, and the represented Laplacian has inverse length
squared. Kinetic scaling then converts that unit to energy. A bare floating
array cannot communicate this path on its own.

The report should also retain negative observations. If a coarse grid breaks an
expected asymptotic trend, that point is not improved by deletion after the
fact. It may reveal that the sequence has not entered the asymptotic regime.
The appropriate conclusion is bounded to the tested sequence.

\section{Common mistakes}

\begin{enumerate}
  \item Treating fractional coordinates as Cartesian coordinates when the cell
        is skew.
  \item Using only diagonal entries of $G^{-1}$ for a nonorthogonal basis.
  \item Confusing a $3\times2$ embedded plane with a $3\times3$ slab.
  \item Applying a seam phase for only one axis when a corner wraps several
        axes.
  \item Adding a potential matrix whose ordering or units differ from the
        Laplacian representation.
  \item Calling $L_h$ a kinetic-energy matrix before supplying the sign,
        coefficient, and mass model.
  \item Interpreting Hermiticity or a passing software test as convergence.
  \item Interpreting convergence to the stated synthetic operator as evidence
        that the operator models a real material adequately.
\end{enumerate}

\section{Exercises}

\begin{enumerate}
  \item Derive Equation~\eqref{eq:training-metric-laplacian} from the chain rule
        for a constant full-rank basis.
  \item Derive the four-corner stencil in
        Equation~\eqref{eq:training-mixed-second-difference} by composing two
        centered first differences.
  \item For the skew basis in
        Equation~\eqref{eq:training-skew-basis}, write the complete
        two-dimensional discrete stencil coefficients in terms of $a$,
        $\theta$, $h_1$, and $h_2$.
  \item Show that the zero-twist constant mode is annihilated by every pure and
        mixed stencil contribution.
  \item Determine the seam quotient $\mathbf q$ and phase for a corner
        displacement from $(N_1-1,N_2-1)$ to
        $(N_1,N_2)$.
  \item Explain why a scalar potential sampled at grid points is diagonal in
        this real-space representation but need not be diagonal after transport
        to a plane-wave representation.
  \item Propose independent oracles for shape/order, Hermiticity, the orthogonal
        reduction, the skew-cell symbol, second-order refinement, and unit
        propagation. Classify each as software or numerical verification.
  \item State the additional physical information required to replace
        Equation~\eqref{eq:training-isotropic-kinetic} by an anisotropic
        multivalley silicon effective-mass operator.
\end{enumerate}

\section{Evidence boundary and next use}

The software dependency and consuming-project smoke test establish that the
accepted public records compose through the embedded-plane, intrinsic-metric,
strict-serialization, and scalar-Laplacian route. PhysKit's own synthetic tests
cover a broader numerical surface. Neither set of tests is a new calculation
for Appendix~\ref{app:two-dimensional-reduction}.

The planned M4 calculation will be the first monograph-owned opportunity to use
the maintained nonorthogonal parent route in a prospectively frozen controlled
comparison. Until its protocol, retained result, independent reconstruction,
and manifest pass, this chapter remains training material and software
capability documentation rather than calculated numerical evidence.
